\documentclass[12pt]{article}
\usepackage[utf8]{inputenc}
\usepackage[T1]{fontenc}
\usepackage{hyperref}
\title{North Tairora Language Resources for K-12 Curriculum Development}
\author{Compiled by Claude Sonnet 5 (claude-sonnet-5) for Andr\'as Kornai}
\date{}
\begin{document}
\maketitle

\section{Introduction}
North Tairora (ISO 639-3 tbg) is a Kainantu-branch language of Eastern Highlands Province, Papua New Guinea. It is spoken south and southeast of Kainantu. The New Guinea World entry gives about 5,000 speakers (2010). Youth are reported to be increasingly using Tok Pisin.

\section{Primary Educational Resources}
No school materials were found.

\section{Linguistic and Reference Materials}
\begin{itemize}
\item New Guinea World entry for North Tairora. \url{https://newguineaworld.linguistik.uzh.ch/families/trans-new-guinea/eastern-highlands-kratke-range/eastern-highlands/kainantu/east-kainantu/tairora/north-tairora}
\item Wikipedia article on the Tairora language. It reportedly cites the Vincent and Vincent grammars (1964) and a Tairora--English dictionary by Vincent and Kaave (2010). I could not confirm the dictionary's availability or whether it covers North Tairora specifically. \url{https://en.wikipedia.org/wiki/Tairora_language}
\item Numerals data page for North Tairora. \url{https://lingweb.eva.mpg.de/channumerals/Tairora-North.htm}
\item Omniglot page on Tairora. \url{https://mail.omniglot.com/writing/tairoa.htm}
\end{itemize}

\section{Religious and Community Resources}
\begin{itemize}
\item Joshua Project page for Northern Tairora. It is general information and a weaker source. \url{https://joshuaproject.net/index.php/people_groups/15197}
\end{itemize}

\section{Digital and Audio Resources}
None were found.

\section{Curriculum Development Context}
Descriptive work exists, but no educational materials were found. The next step would be to locate the Vincent grammars and the 2010 dictionary, and to check orthography and Bible translation status with SIL PNG and local churches.
\end{document}